\subsection{Canonicalizing $\R^{d \times n}$ with respect to $\O(d)$ and $\SO(d)$}\label{sec:canonicalizing-matrices-rotations}
\begin{definition}
    The \emph{orthogonal group} of degree $n$ is the set
    \[
        \O(n) := \set{Q \in \R^{n \times n} : Q Q^T = I}.
    \]
\end{definition}

\begin{definition}
    The \emph{special orthogonal group} of degree $n$ is the set
    \[
        \SO(n) := \set{Q \in \R^{n \times n} : \det(Q) = 1}.
    \]
\end{definition}

\begin{remark}
    Observe that $\SO(n) \sub \O(n)$.
\end{remark}

\begin{definition}
    The set of \emph{real symmetric matrices of degree $n$} is
    \[
        \S^n := \set{A \in \R^{n \times n} : A = A^T}.
    \]
\end{definition}

\begin{definition}
    The set of \emph{positive semi-definite matrices of degree $n$} is
    \[
        \S^n_+ := \set{A \in \S^n : \forall x \in \R^n, x^T A x \geq 0}.
    \]
\end{definition}

\begin{definition}
    The set of \emph{negative semi-definite matrices of degree $n$} is
    \[
        \S^n_- := \set{A \in \S^n : \forall x \in \R^n, x^T A x \leq 0}.
    \]
\end{definition}

\begin{definition}
    The set of \emph{positive semi-definite matrices of degree $n$ with at most rank $d$} is
    \[
        \S^n_{\leq d} := \set{A \in \S^n_+ : \rank(A) \leq d}.
    \]
\end{definition}

\begin{remark}
    $\S^n_{\leq d} = \S^n_{\leq n}$ for any $d \geq n$.
\end{remark}

\begin{proposition}
    $A \in \S^n_+$ if and only if there is some matrix $B \in \R^{n \times n}$ such that $A = B^T B$ \cite{wikipedia:definite-matrices}.
\end{proposition}

\begin{corollary}
    $\S^n_+ = \S^n_{\leq n}$.
\end{corollary}

\begin{definition}
    The set of \emph{positive semi-definite matrices of degree $n$ with rank $d$} is
    \[
        \S^n_d := \set{A \in \S^n_+ : \rank(A) = d}.
    \]
\end{definition}

\begin{proposition}\label{prop:psd-iff-transpose-product}
    $A \in \S^n_d$ if and only if there is some matrix $B \in \R^{d \times n}$ such that $A = B^T B$ with $\rank(B) = d$. \cite{wikipedia:definite-matrices}.
\end{proposition}

\begin{lemma}\label{lem:Od-cong-psd}
    With $O(d)$ acting on $\R^{d \times n}$ via matrix multiplication on the left, $\R^{d \times n} / \O(d)$ is homeomorphic to the space $S := \set{A^T A : A \in \R^{d \times n}}$.
    \begin{proof}

        To begin, consider the map $\phi : \R^{d \times n}$ given by $\phi(A) = A^T A$.
        Now let $A,B \in \R^{d \times n}$ such that $[A] = [B] \in \R^{d \times n} / O(d)$, i.e. there exists $Q \in O(d)$ such that $A = QB$.
        This means that $A^T A = (QB)^T QB = B^T Q^T Q B = B^T B$, so $\phi$ induces a map $\Hat{\phi} : \R^{d \times n} / O(d) \to S$ via $\Hat{\phi}([A]) = A^T A$.
        Note that $\Hat{\phi}$ is bijective by construction, with inverse $\inv{\Hat{\phi}} : S \to \R^{d \times n} / O(d)$ given by $\inv{\Hat{\phi}}(A^T A) = [A]$.
        Both of these maps are continuous, so we are done.
    \end{proof}
\end{lemma}

\begin{corollary}\label{cor:Rnd/Od-cong-Snd}
    $\R^{d \times n} / O(d) \cong \S^n_{\leq d}$.
    \begin{proof}
        This follows directly from proposition \ref{prop:psd-iff-transpose-product} and lemma \ref{lem:Od-cong-psd}.
    \end{proof}
\end{corollary}

With corollary \ref{cor:Rnd/Od-cong-Snd} in hand, we have that any invariant map $f : \R^{d \times n} \to Y$ factors through $\S^n_d$ and so it suffices to reduce to consider maps from $\S^n_d$ to $Y$.

\begin{definition}
    Given a group $G$ and a $G$-set $X$, the \emph{free locus} of $X$ is
    \[
        X_{free}
        :=
        \set{x \in X : \forall g \in G \sminus \set{e}, g x \neq x}
        =
        \set{x \in X : G_x = \set{e}}
    \]
\end{definition}

\begin{lemma}\label{lem:free-iff-full-rank}
    With $O(d)$ acting on $\R^{d \times n}$ via matrix multiplication, $\R^{d \times n}_{free} = \R^{d \times n}_*$.
    \begin{proof}
        First, let $A \in \R^{d \times n}$ such that $\rank(A) = k < d$.
        Then there are some orthonormal bases $\set{e_1,...,e_k}$ of $\col(A)$ and $\set{e_{k+1},...,e_d}$ of $\nul(A)$.
        Let $B = [e_1 \cdots e_d]^T \in O(d)$ be the change of basis matrix from the standard basis for $\R^d$ to $\set{e_1,...,e_d}$.
        Then $A$ can be written in the basis $\set{e_1,...,e_d}$ as a matrix
        \[
            \Hat{A} = 
            BA = 
            \begin{bmatrix}
                \Hat{a}_{1,1} & \cdots & \Hat{a}_{n,1} \\
                \vdots & & \vdots \\
                \Hat{a}_{k,1} & \cdots & \Hat{a}_{n,k} \\
                0 & \cdots & 0 \\
                \vdots & & \vdots \\
                0 & \cdots & 0
            \end{bmatrix}.
        \]
    Now define the matrix $R \in \R^{d \times d}$ by $R_{i,i} = 1$ if $i \leq k$, $R_{i,i} = -1$ if $i > k$, and $R_{i,j} = 0$ if $i \neq j$ ($R$ simply flips the sign of the last $d-k$ coordinates) and observe that clearly $R \in O(d)$.
    Define $Q = B^T R B$, and then $Q \in O(d)$ since
    \[
        Q^T Q
        = (B^T R B)^T B^T R B
        = B^T R^T B B^T R B
        = I.
    \]
    By construction we have that $R\Hat{A} = \Hat{A}$, so it follows that $QA = A$.
    However, without loss of generality $Q \neq I$, so we have that $A \not\in \R^{d \times n}_{free}$.
    
    Conversely, let $A \in \R^{d \times n}_*$ and let $Q \in O(d)$ such that $QA = A$.
    Since $\col(A) = \R^d$, this means that $Q$ fixes all of $\R^d$ and hence $Q = I$.
    Therefore $A \in \R^{d \times n}_{free}$.
    \end{proof}
\end{lemma}

\begin{corollary}
    $O(d)$ does not act freely on any subset of $\R^{d \times n}$ if $n < d$.
\end{corollary}

\begin{corollary}
    $\R^{d \times n}_* / O(d) \cong \S^n_d$.
    \begin{proof}
        This follows from proposition \ref{prop:psd-iff-transpose-product}, lemma \ref{lem:Od-cong-psd}, and lemma \ref{lem:free-iff-full-rank}.
    \end{proof}
\end{corollary}

We now present a somewhat simplified proof of the result of \cite{colman2013equivariant}  that $(\R^{d \times n}, O(d))$ and $(\R^{d \times n}, SO(d))$ lack continuous canonicalizations for $n > d$ and $n \geq d$ respectively.
Propositions \ref{prop:PSD-mod-scale-cong-ball} and \ref{prop:mod-SOn-cong-PSD} and theorem \ref{thm:Casson} are due to Casson, as presented in theorem 1.1 of \cite{kendall1999shape}.

\begin{definition}
    For a finite dimensional matrix $A \in \R^{n \times m}$, the \emph{spectrum} of $A$ is the set of (possibly complex) eigenvalues of $A$, denoted $\sigma(A)$.

    Note: this is a slight abuse of the standard definition of the spectrum of a linear operator, but since we are in finite dimensions this definition coincides with the more general one.
\end{definition}

\begin{lemma}\label{lem:sym-tr0-eigenbound}
    Let $E \in \S^n \sminus \set{0}$ with $\tr(E) = 0$, and let $\eps \in \R$ be the least eigenvalue of $E$.
    Then $\abs{\eps} > \f{\norm{E}}{n}$ where $\norm{\cdot}$ denotes the operator norm.
    \begin{proof}
        \Afsoc that $\eps \leq \f{\norm{E}}{n}$.
        Let $\set{\lambda_i}$ and $\set{\mu_j}$ be the positive and negative eigenvalues of $E$ respectively.
        Recall that $\tr(E) = \sum_i \lambda_i + \sum_j \mu_j$, so $E$ must have at least one positive and one negative eigenvalue since $E \neq 0$.
        Also note that we must have that $\abs{\mu_j} \leq \abs{\eps}$ for all $j$.
        From functional analysis we know that $\norm{E} = \max\set{\abs{\min \sigma(E)}, \abs{\max \sigma(E)}}$, so since $\abs{\eps} < \norm{E}$ we must have some positive eigenvalue $\lambda_j = \norm{E}$.
        Then we have
        \[
            \norm{E}
            \leq \sum_i \lambda_i
            = \sum_j \abs{\mu_j}
            \leq (n-1) \abs{\eps}
            < n \abs{\eps}
            \leq \norm{E}
        \]
        which is a contradiction.
        Therefore we conclude that $\abs{\eps} > \f{\norm{E}}{n}$.
    \end{proof}
\end{lemma}

\begin{proposition}\label{prop:PSD-mod-scale-cong-ball}
    Let $Y_+ = \S^n_+ \setminus \set{0}$ be the set of non-zero positive semidefinite matrices of degree $n$, let $Y_0 = \set{A \in Y_+ : \det(A) = 0}$, and let $\Rp$ be the group of positive real numbers under multiplication.
    If $\Rp$ acts on the left of $Y_+$ via scalar multiplication, then $Y_+/\Rp$ is homeomorphic to a closed ball of dimension $n(n+1)/2 - 1$.
    Additionally, $Y_0/\Rp$ is homeomorphic to the ball's boundary.
    
    \begin{proof}
        The basic idea of this proof is that we observe that $Y_+$ is a convex cone without its point, so we simply take a slice of it and homeomorphically map that slice to a disk.
        Then it is reasonably clear that such a slice is exactly the cone quotient by scaling, so we conclude there.
        The difficulty lies in showing the slice is actually a topological disk, the details of which lie in this proof.

        Let $\angles{A,B} := \tr(AB)$ be the Frobenius inner product on $\S^n$, which induces the Frobenius norm $\norm{\cdot}_F$.
        Then let $I^\perp$ be orthogonal complement of $I$ in $\S^n$, let $D: = (I + I^\perp) \cap Y_+$, and let $B := \clo{B(0,1)} \cap I^\perp$ be the closed unit ball in $I^\perp$ under the operator, which has dimension $\dim(I^\perp) = \dim(Y_+) - \dim(\Span\set{I}) = n(n+1)/2 - 1$.
        Observe that $D$ is homeomorphic to $Y_+ / \Rp$ via the quotient map $E \mapsto [E]$ with continuous inverse $[E] \mapsto \f{1}{\norm{\proj_I E}} E$.

        Define the function $\mu : \S^n \to \R$ by $\mu(A) = \min \sigma(A)$.
        $\mu$ is well-defined since real symmetric matrices have finitely many eigenvalues, all of which are real, and it follows from Weyl's inequality that $\mu$ is continuous.
        For any $A \in \S^n$, observe that $I + A \in Y_+$ if and only if $\mu(I + A) \geq 0$.

        We now define $\Phi : B \to D$ via $\Phi(E) =  I - \f{\norm{E}}{\mu(E)}E$ if $E \neq 0$ and $\Phi(0) = I$.
        We will show that $\Phi$ is the desired homeomorphism.
        Observe that $\Phi(E) \in (I + I^\perp)$ by construction, but it remains to show that $\Phi(E) \in Y_+$.
        This is trivially true for $E = 0$, so we consider $E \neq 0$.
        Note that $E \in E^\perp$ means that $\tr(E) = 0$, meaning that the sum of $E$'s eigenvalues is $0$, so $E$ must have at least one negative eigenvalue since $E \neq 0$.
        Hence, $\mu(E) < 0$.
        Then we see that $E \in Y_+$ since
        \begin{align*}
            \min_{\norm{x} = 1} x^T \Phi(E) x
            &= \min_{\norm{x} = 1} x^T \bigp{I - \f{\norm{E}}{\mu(E)}E} x \\
            &= 1  - \f{\norm{E}}{\mu(E)} \min_{\norm{x} = 1} x^T E x \\
            &= 1 - \norm{E} \\
            &\geq 0.
        \end{align*}
        This also shows that $\Phi(E) \in Y_0$ if and only if $\norm{E} = 1$.
        The continuity of $\Phi$ at every point other than $0$ is evident from its construction.
        To see $\Phi$'s continuity at $0$, first observe that
        \[
            \f{\mu(E)}{\norm{E}}
            = \f{\min_{\norm{x} = 1} x^T E x}{\norm{E}}
            = \min_{\norm{x} = 1} x^T \f{E}{\norm{E}} x
            = \mu\bigp{\f{E}{\norm{E}}}
        \]
        so
        \[
            \Phi(E) = I - \f{1}{\mu\bigp{\f{E}{\norm{E}}}} E.
        \]
        Note that $\Phi(E)$ is continuous if and only if $\f{1}{\mu\bigp{\f{E}{\norm{E}}}} E$ is, and since $\norm{\f{E}{\norm{E}}} = 1$ we use lemma \ref{lem:sym-tr0-eigenbound} to get that $\mu\bigp{\f{E}{\norm{E}}} > \f{1}{n}$ and  compute
        \[
            \lim_{E \to 0} \norm{\f{1}{\mu\bigp{\f{E}{\norm{E}}}} E}
            < \lim_{E \to 0} \norm{n E}
            = 0
        \]
        so $\lim_{E \to 0} \Phi(E) = I$ and $\Phi$ is continuous.
        
        We show that $\Phi$ is indeed a homeomorphism by constructing its inverse $\inv{\Phi} : D \to B$ via $\inv{\Phi}(I + E) = -\f{\mu(E)}{\norm{E}} E$ for $E \neq 0$ and $\inv{\Phi}(I) = 0$.
        To see that this is indeed the inverse of $\Phi$, we first compute
        \begin{align*}
            \Phi(\inv{\Phi}(I + E))
            &= \Phi\bigp{-\f{\mu(E)}{\norm{E}}E} \\
            &= I - \f{\norm{-\f{\mu(E)}{\norm{E}}E}}{\mu\bigp{-\f{\mu(E)}{\norm{E}}E}}\bigp{-\f{\mu(E)}{\norm{E}} E} \\
            &= I + \f{-\f{\mu(E)}{\norm{E}}\norm{E}}{-\f{\mu(E)}{\norm{E}}\mu\bigp{E}}\f{\mu(E)}{\norm{E}} E \\
            &= I + E
        \end{align*}
        when $E \neq 0$, and the computation is trivial when $E = 0$.
        Similarly, we have
        \begin{align*}
            \inv{\Phi}(\Phi(E))
            &= \inv{\Phi}\bigp{I - \f{\norm{E}}{\mu(E)} E} \\
            &= -\mu\bigp{\f{-\f{\norm{E}}{\mu(E)} E}{\norm{-\f{\norm{E}}{\mu(E)} E}}}\bigp{-\f{\norm{E}}{\mu(E)} E} \\
            &= \f{-\f{\norm{E}}{\mu(E)}}{-\f{\norm{E}}{\mu(E)} \norm{E}} \mu(E) \f{\norm{E}}{\mu(E)} E \\
            &= E.
        \end{align*}
        Thus we have that $\inv{\Phi}$ is indeed the inverse of $\Phi$.
        Lastly, $\inv{\Phi}$ is again clearly continuous everywhere except $0$ by construction, and at $0$ we have by the continuity of $\mu$ that
        \[
            \lim_{E \to 0} \norm{\inv{\Phi}(E)}
            = \lim_{E \to 0} \norm{-\f{\mu(E)}{\norm{E}} E}
            = \lim_{E \to 0} \abs{\mu(E)} \f{\norm{E}}{\norm{E}}
            = 0.
        \]
        We conclude that $\Phi$ is a homeomorphism, hence $Y_+/\Rp \cong B$ with $Y_0/\Rp \cong \bdry B$.
    \end{proof}
    \begin{proof}[Alternative proof]
        First, we wish to show that the identity matrix $I$ is interior point of $Y_+ \sub \S^n$ in the subspace topology.
        Recall all norms are equivalent on a finite dimensional vector space, so it suffices to find an open ball under the operator norm.
        Define the function $\mu : \S^n \to \R$ by $\mu(A) = \min \sigma(A)$.
        $\mu$ is well-defined since real symmetric matrices have finitely many eigenvalues, all of which are real, and it follows from Weyl's inequality that $\mu$ is continuous.
        For any $A \in \S^n$, observe that $I + A \in Y_+$ if and only if $\mu(I + A) \geq 0$.
        By the continuity of $\mu$, there is some $\delta$ such that $\norm{A} < \delta$ implies $\abs{\mu(I) - \mu(I + A)} < 1$, meaning $\mu(I+A) > 0$ since $\mu(I) = 1$.
        Therefore $B(I,\delta) \sub Y_+$, as desired, so $\dim Y_+ = \dim \S^n = n(n+1)/2$.
        
        Next, $Y_+$ is clearly convex since the sum of any two positive semidefinite matrices is also positive semidefinite.
        % TODO: formalize this
        From here, one can simply consider the intersection $Y_+ \cap (I + I^\perp)$ which is a bounded, convex subset of $Y_+$ and is hence homeomorphic to the ball of dimension $n(n+1)/2 - 1$.
        Additionally, $\bdry Y_+ = Y_0$ since any singular matrix can be perturbed to have a negative eigenvalue, so $Y_0 \cap (I + I^\perp)$ is homeomorphic to the ball's boundary.
    \end{proof}
\end{proposition}

\begin{proposition}\label{prop:mod-SOn-cong-PSD}
    Let  $X_+ = \set{A \in \R^{n \times n} \sminus \set{0} : \det(A) \geq 0}$ with left $\SO(n)$ action given by matrix multiplication, and let $Y_+ = \S^n_+ \sminus \set{0}$.
    Additionally, let $X_0 \sub X_+$ and $Y_0 \sub Y_+$ be the subspaces of singular matrices.
    Then $X_+ / \SO(n) \cong Y_+$ with $X_0 / \SO(n) \cong Y_0$.
    \begin{proof}
        This follows from a nearly identical argument to that of lemma \ref{lem:Od-cong-psd} and corollary \ref{cor:Rnd/Od-cong-Snd}, simply with the added care to only use matrices with nonnegative determinant.
    \end{proof}
\end{proposition}

\begin{theorem}[Casson]\label{thm:Casson}
    Let $X = \R^{n \times n} \sminus \set{0}$ with $\SO(n)$ and $\Rp$ acting on the left via matrix and scalar multiplication respectively, and let $S$ be the sphere of dimension $n(n+1)/2 - 1$.
    Then $X / (\Rp \times \SO(n)) \cong S$.
    \begin{proof}
        First observe that the actions of $\Rp$ and $\SO(n)$ on $X$ commute, so we have that
        $X / (\Rp \times \SO(n)) \cong (X / \SO(n)) / \Rp$.
        Let $X_0,X_+,X_-$ be the subspaces of $X$ of matrices with zero, nonnegative, and nonpositive determinant respectively, and observe that clearly $X \cong X_+ \cup_{X_0} X_-$.
        Observe that any $\SO(n)$ orbit is entirely contained in $X_0$, $X_+$, or $X_-$, since $\SO(n)$-actions do not change the determinant.
        Additionally, $X_+ \cong X_-$ via your favorite map that flips the sign of the determinant (e.g. multiplying the last column by $-1$).
        It then follows that
        \[
            (X / \SO(n)) / \Rp
            \cong (X_+ / \SO(n)) / \Rp
            \cup_{(X_0 / \SO(n)) / \Rp}
             (X_+ / \SO(n)) / \Rp.
        \]
        By propositions \ref{prop:PSD-mod-scale-cong-ball} and \ref{prop:mod-SOn-cong-PSD} we know that $(X_+ / \SO(n))/\Rp$ and $(X_0 / \SO(n))/\Rp$ are homeomorphic to the ball of dimension $n(n+1)/2 - 1$ and its boundary respectively, so we conclude that $X / (\Rp \times \SO(n))$ is homeomorphic to the sphere of dimension $n(n+1)/2 - 1$.
    \end{proof}
\end{theorem}

With Casson's theorem in hand we can now rather quickly prove two of \cite{dym2024equivariant}'s impossibility results.

\begin{lemma}\label{lem:induce-canon-truncate}
    (From \cite{dym2024equivariant})
    Let $G \sub \R^{d \times d}$ be a group acting on $\R^{d \times n}$ on the left via matrix multiplication.
    Then a canonicalization for $\R^{d \times n}$ induces a canonicalization for $\R^{d \times m}$ for any $m \leq n$.
    \begin{proof}
        Let $f : \R^{d \times n} \to \R^{d \times n}$ be a canonicalization, and consider the map $\Hat{f} : \R^{d \times m} \to \R^{d \times m}$ given by $\Hat{f}(A) = f([A \mid 0])[:,:m]$ ($[A \mid 0]$ indicates padding $A$ on the right with zeroes so $[A \mid 0] \in \R^{d \times n}$).
        Then for any $Q \in G$ and $A \in \R^{d \times m}$,
        \[
            \Hat{f}(QA)
            = f(Q[A \mid 0])[:,:m]
            = (\Hat{Q}[A \mid 0])[:,:m]
            = \Hat{Q}A
        \]
    for some $\Hat{Q} \in G$.
    Note that $\Hat{f}(QB) = \Hat{Q}A$ any $Q \in G$, and $B \in GA$ since $f$ is a canonicalization, so we conclude that $\Hat{f}$ is a canonicalization as well.
    \end{proof}
\end{lemma}

\begin{theorem}\label{thm:no-SOd-canon-for-dxn}
    If $n \geq d$, $\R^{d \times n}$ has no canonicalization with respect to $\SO(d)$.
    \begin{proof}
        By lemma \ref{lem:induce-canon-truncate} we know that any canonicalization of $\R^{d \times n}$ induces a canonicalization for $\R^{d \times d}$, so we reduce to considering this case.
        Now, \afsoc that we have a canonicalization for $\R^{d \times d}$, and recall that this induces a section  of the quotient map $\R^{d \times d} \to \R^{d \times d}/\SO(d)$.
        Observe that such a section also induces a section on any $\SO(d)$-invariant subspace of $\R^{d \times d}$, so we specifically consider the unit sphere $S^{d^2-1} \sub \R^{d \times d}$ under the Frobenius norm.
        It is easy to see that $S^{d^2-1}$ is invariant under $\SO(d)$ transformations since orthogonal transformations do not change the Frobenius norm, and it is also clear that $S^{d^2-1}$ is homeomorphic to the quotient $(\R^{d \times d} \sminus \set{0}) / \Rp$ since one can always scale a non-zero matrix to have norm $1$.
        Hence, we have a section $s : S^{d^2-1}/\SO(d) \to S^{d^2-1}$ of the quotient map $q : S^{d^2-1} \to S^{d^2-1}/\SO(d)$.

        Recall by theorem \ref{thm:Casson} that the quotient $S^{d^2-1}/\SO(d)$ is homeomorphic to the sphere of dimension $d(d+1)/2-1$, $S^{\f{d(d+1)}{2} - 1}$.
        So, up to slight abuse of notation, we have induced maps on homology
        \[
            H_k(S^{\f{d(d+1)}{2} - 1})
            \tto{s_\sharp}
            H_k(S^{d^2 - 1})
            \tto{q_\sharp}
            H_k(S^{\f{d(d+1)}{2} - 1})
        \]
        such that $q_\sharp \circ s_\sharp$ is still the identity.
        However, when $k = d(d+1)/2-1$ this is impossible, since in this case $H_k(S^{\f{d(d+1)}{2} - 1}) \cong \Z$ and $H_k(S^{d^2-1}) \cong 0$.
        Therefore we conclude that there can be no canonicalization of $\R^{d \times n}$ with respect to $\SO(d)$ for $n \geq d$.
    \end{proof}
\end{theorem}

\begin{theorem}\label{thm:no-Od-canon-for-dxn}
    If $n > d$, $\R^{d \times n}$ has no canonicalization with respect to $\O(d)$.
    \begin{proof}[Proof 1]
        This proof proceeds almost identically to that of theorem \ref{thm:no-SOd-canon-for-dxn}.
        We first reduce to considering $\R^{d \times (d+1)}$, then find $S^{d(d+1) - 1} \sub \R^{d \times (d+1)}$, which is homeomorphic to $\R^{d \times (d+1)} / \Rp$.
        Next, observe that $S^{d(d+1) - 1}/O(d)$ is homeomorphic to $\S^{d+1}_{\leq d}/\Rp$ by corollary \ref{cor:Rnd/Od-cong-Snd}, which is in turn homeomorphic to the boundary of the ball of dimension $(d+1)(d+2)/2 - 1$, i.e. the sphere of one dimension less than that, by proposition \ref{prop:PSD-mod-scale-cong-ball}.
        Then a contradiction once again arises by considering the induced maps on homology.
        We require $d \geq 4$ here so that $d^2 - 1 > (d+1)(d+2)/2 - 2$ and this last step actually produces a contradiction.
        For lower values of $d$, one can use \cite{dym2024equivariant}'s proof.
    \end{proof}
    This result also has a more direct proof.
    \begin{proof}[Proof 2]
        \Afsoc that we have some canonicalization of $\R^{d \times n}$ and the associated section $s : \S^n_{\leq d} \to \R^{d \times n}$.
    \end{proof}
\end{theorem}

\begin{proposition}\label{prop:SOd-secat-2}
    Let $\SO(d)$ act on the left of $\R^{d \times d}_*$ via matrix multiplication and let $q : \R^{d \times d}_* \to \R^{d \times d}_* / \SO(d)$ be the quotient map.
    Then $\ssecat(q) = 2$.
    \begin{proof}
        First, observe that by theorem \ref{thm:no-SOd-canon-for-dxn}, proposition \ref{prop:canon-is-section}, remark \ref{rem:secat1-iff-section}, and lemma \ref{lem:subsecat-leq-secat} we have that $\ssecat(q) > 1$.
        Hence, it remains to show that $\ssecat(q) \leq 2$.
        By proposition \ref{prop:psd-iff-transpose-product} we know that $\R^{d \times d}$ is homeomorphic to the (disjoint) union of full rank matrices in $\S^d_+$ and $\S^d_-$ (here denoted $S_+$ and $S_-$) via the map $q : \R^{d \times d}_* \to S_+ \sqcup S_-$ defined as follows:
        \[
            q(A) = 
            \begin{cases}
                A^T A &\text{ if } \det(A) > 0 \\
                -A^T A &\text{ if } \det(A) < 0
            \end{cases}.
        \]
        Observe that this well-defined since $A^T A$ is always positive semi-definite with $\rank(A^T A) = \rank(A)$, so $A^T A \in S_+$ and $-A^T A \in S_-$.
        Additionally, $q$ is clearly continuous.

        Next, we define $s_+ : S_+ \to \R^{d \times d}$ by $s_+(A) = A^{\f12}$ and $s_- : S_- \to \R^{d \times d}$ by $s_-(A) = T(-A)^{\f12}$, where $M^{\f12}$ denotes the matrix square root where $T$ is some orthogonal transformation that flips the sign of its input's determinant, i.e. $T \in O(d) \sminus SO(d)$.
        Both of these maps are clearly well-defined and continuous by construction.
        For $A \in S_+$, we then compute
        \[
            q(s_+(A))
            = q(A^{\f12})
            = (A^{\f12})^T (A^{\f12})
            = A
        \]
        and similarly for $A \in S_-$ we have
        \[
            q(s_-(A))
            = q(T(-A)^{\f12})
            = -(T(-A)^{\f12})^T (T(-A)^{\f12})
            = -(-A)^{\f12} T^T T (-A)^{\f12}
            = A.
        \]
        Therefore both $s_+$ and $s_-$ are local sections of $q$ one $S_+$ and $S_-$ respectively, and both $S_+$ and $S_-$ are open in the quotient by construction so we conclude that $\ssecat(q) \leq 2$.
    \end{proof}
\end{proposition}

We would also like to compute or at least bound the sectional category of the quotients $\R^{d \times n} \to \R^{d \times n} / \SO(d)$ and $\R^{d \times n} / \O(d)$ in general.
To do this, it is useful to consider the Stiefel and Grassmann manifolds as in section \ref{sec:bg-stiefel-whitney-classes}.


\begin{remark}
    (From \cite{bendokat2024grassmanhandbook})
    The Grassmann manifold $\Gr(n,k)$ can also be realized as a quotient of the Stiefel manifold $\St(n,k)$ by relating matrices that are basis for the same subspace, i.e. those that differ by an orthogonal transformation.
    This is captured by the action of $O(k)$ on $\St(n,k)$ by right multiplication, hence $\Gr(n,k) \cong \St(n,k) / O(k)$.
\end{remark}

Both of the the Stiefel and Grassman manifolds are compact, and have corresponding non-compact variants:

\begin{definition}
    The \emph{non-compact Stiefel manifold} $\St'(n,k)$ is the set of basis matrices for $k$-dimensional subspaces of $\R^n$, which can be represented via
    \[
        \St'(n,k) := \set{U \in \R^{n \times k} : \rank(U) = k}
    \]
\end{definition}

\begin{definition}
    The \emph{non-compact Grassmann manifold} $\Gr'(n,k)$ is $\St'(n,k) / O(k)$ where $O(k)$ acts on the right of $\St'(n,k)$ via matrix multiplication.
\end{definition}

\begin{remark}\label{rem:full-rank-hom-stiefel}
    The transposition map $A \mapsto A^T$ gives a homeomorphism between $\R^{d \times n}_*$ and $\St'(n,d)$, hence $\R^{d \times n}_* / O(d) \cong \Gr'(n,d)$.
\end{remark}

In view of remark \ref{rem:full-rank-hom-stiefel}, we now wish to compute/bound the sectional category of the projection $\pi : \St(n,k) \to \Gr(n,k)$, treating $\St(n,k)$ as the unit sphere in $\R^{d \times n}_*$.
But recall that by corollary \ref{cor:nil-st-gr-proj-eq-cup-gr} this can be done by computing/bounding the cup length of $\Gr(n,k)$.

\begin{proposition}\label{prop:cup-Grnk-monotonicity}
    (From \cite{pavesic2021} following proposition 3.5)
    For $n \leq m$, $\cupln(\Gr(n,k)) \leq \cupln(\Gr(m,k))$.
\end{proposition}

\begin{proposition}\label{prop:cup-Grn2}
    (From \cite{pavesic2021}, proposition 3.3)
    If $2^s < n \leq 2^{s+1}$ for some integer $s$, then
    \[
        \cupln(\Gr(n,2))
        = \begin{cases}
            2^{s+1} - 2 &\tif n = 2^s + 1 \\
            2^{s+1} - 1 &\totherwise.
        \end{cases}
    \]
\end{proposition}

\begin{proposition}\label{prop:cup-Grn3}
    (From \cite{stong1982}, p. 104)
    If $2^s < n \leq 2^{s+1}$ for some positive integer $s$, we have
    \[
        \cupln(\Gr(n,3))
        = \begin{cases}
            2^{s + 2} - 3 \cdot 2^{p - 1} - 4 &\text{ if } n = 2^{s+1} - 2^{p} + 1 \text{ for some } p \geq 1 \\
            2^{s + 2} - 3 \cdot 2^{p - 1} - 4 + t &\text{ if } n = 2^{s+1} - 2^{p} + t\\
            &\quad\text{ for some } p \geq 1, 2 \leq t \leq 2^{p-1} \\
            2^{s + 2} - 5 &\text{ if } n = 2^{s+1} \\
        \end{cases}
    \]
    From this, we have a tight lower bound of $\cupln(\Gr(n,3)) \geq 2n - 5$.
\end{proposition}

\begin{proposition}\label{prop:cup-Grn4}
    Letting $s$ range over the positive integers, we have
    \[
        \cupln(\Gr(n,4))
        = \begin{cases}
            3 \cdot 2^s - 7 + 2^{r+1} + t &\tif n = 2^s + 1 + 2^r + t \\
            &\quad\text{ for some } 0 \leq r < s, 0 \leq t < 2^r \\
            3 \cdot 2^s - 7 &\tif n = 2^s + 1
        \end{cases}
    \]
    From this, we have a tight lower bound of $\cupln(\Gr(n,4)) \geq \f{9}{4} n - 8$.
\end{proposition}

\begin{proposition}\label{prop:cup-Grn5-partial}
    (From \cite{stong1982}, p.104)
    For $s \geq 3$, $\cupln(\Gr(2^s+1,5)) = 3 \cdot 2^s + 2^{s-2} - 9$.
\end{proposition}

\begin{corollary}
    Combining propositions \ref{prop:cup-Grn5-partial} and \ref{prop:cup-Grnk-monotonicity}, we have that
    \[
        \cupln(\Gr(n,5)) \geq \f{13}{8} n - \f{85}{8}
    \]
\end{corollary}

For more general $k$, we have the following bound.

\begin{proposition}\label{prop:cup-Grnk-bound}
    (Corollary of \cite{stong1982}, p. 103)
    If $2 \leq k \leq n/2$ and $2^s < n \leq s^{s+1}$ for some integer $s$, then
    \[
        \cupln(\Gr(n,k))
        \geq \begin{cases}
            2^{s+1} - 2 &\tif n = 2^s+1 \\
            2^{s+1} - 1 &\totherwise.
        \end{cases}
    \]
    In turn, this means that $\cupln(\Gr(n,k)) \geq n-1$.
\end{proposition}

All of these cup length bounds can thus be used to compute lower bounds for $\ssecat(q)$ with the quotient map $q :\R^{d \times n} \to \R^{d \times n}/\O(d)$, summarized in the following theorem.

\begin{theorem}\label{ref:Od-secat-lower-bounds}
    Let $n \geq d$ and let $q : \R^{d \times n} \to \R^{d \times n}/\O(d)$ be the quotient map.
    Then
    \[
        \secat(q)
        \geq \begin{cases}
            2n - 5 &\tif d = 3 \\
            \f{9}{4}n - 8 &\tif d = 4 \\
            \f{13}{8}n - \f{85}{8} &\tif d = 5 \\
            n - 1 &\tif 6 \leq d \leq \f{n}{2}.
        \end{cases}
    \]
\end{theorem}


We now turn to the case of $q : \R^{d \times n} \to \R^{d \times n}/\SO(d)$.
This case turns out to me markedly more difficult to study, as we now must consider the \emph{oriented} Grassmann manifolds, which less is known about.

\begin{definition}
    The \emph{oriented Grassmann manifold} $\Tilde{\Gr}(n,k)$ is the space of all oriented $k$-dimensional vector subspaces of $\R^n$.
\end{definition}

\begin{remark}
    $\Tilde{\Gr}(n,k) \simeq \R^{d \times n}_* / \SO(d)$.
\end{remark}

\begin{remark}
    There is a natural vector bundle over $\Tilde{\Gr}(n,k)$ given by
    \[
        \Tilde{E}(n,k) := \set{(P,x) \in \Tilde{\Gr}(n,k) \times \R^n : x \in P}\footnote{``$x \in P$'' is a slight abuse of notation here, since $P$ is not a set but rather an oriented plane. However, the meaning is still the obvious one: that $x$ lies in the (unoriented) plane specified by $P$.}
    \]
    with the projection $p : \Tilde{E}(n,k) \to \Tilde{\Gr}(n,k)$.
\end{remark}

\begin{remark}
    There is a natural projection $\Tilde{\pi} : \St(n,k) \to \Tilde{\Gr}(n,k)$ where $\Tilde{\pi}(B)$ is the plane $\Span(B)$ with orientation given by the order of the basis $B$.
\end{remark}

\begin{proposition}\label{prop:oGr-SW-sub-projection-kernel}
    Let $\Tilde{\pi} : \St(n,k) \to \Tilde{\Gr}(n,k)$ be the natural projection.
    Then the induced map $\Tilde{\pi}^* : H^*(\Tilde{\Gr}(n,k)) \to H^*(\St(n,k))$ is zero on the Stiefel-Whitney classes $w_i(\Tilde{E}(n,k))$.
    \begin{proof}
        This argument is exactly the same as proposition \ref{prop:Stiefel-Grassmann-projection-kernel}, except that we no longer have that $H^*(\Tilde{\Gr}(n,k);\Z_2)$ is generated by Stiefel-Whitney classes.
    \end{proof}
\end{proposition}

\begin{definition}
    The \emph{height} of a cohomology class $x \in H^*(X)$ is
    \[
        ht(x) := \sup\set{m \in \Zp : x^m \neq 0}
    \]
    where $x^m$ denotes $\underbrace{x \smile x \smile \cdots \smile x}_{m \text{ times}}$.
\end{definition}

\begin{remark}\label{rem:nil-ker-pi-star-geq-SW-height}
    In light of proposition \ref{prop:oGr-SW-sub-projection-kernel}, we have that
    \[
        \nil\ker \Tilde{\pi}^* \geq \sup_{x \in H^*(\Tilde{\Gr}(n,k))} ht(x).
    \]
\end{remark}

\begin{theorem}
    Let $n/2 \geq d \geq 3$ and let $q : \R^{d \times n} \to \R^{d \times n}/\SO(d)$ be the quotient map.
    Then
    \[
        \secat(q)
        \geq \begin{cases}
            \f{n - d + 6}{2} &\tif n - d + 3 \geq 7 \text{ is odd} \tand n - d + 3 \neq 9,11 \\
            \f{n - d + 5}{2} &\tif n - d + 3 \geq 7 \text{ is even} \tand n - d + 3 \neq 10,12 \\
            5 &\tif n - d + 3 = 9,10,11,12.
        \end{cases}
    \]
    \begin{proof}
        This follows directly from remark \ref{rem:nil-ker-pi-star-geq-SW-height} and the proof of Proposition B in \cite{korbas2006cuplengthbounds}, where the author computes the heights of some Stiefel-Whitney classes.
    \end{proof}
\end{theorem}


We might also wish to giver an upper bound for the sectional category of the two quotient maps considered above.
Unfortunately, I could not find a way to do this.

% TODO: put some graphs of the bounds

